\section{Background and Related Work}
\label{sec:background}

\subsection{Code Representations for Learned Detectors}
Learned vulnerability detection has moved through three representations. Early
sequence models treated a function as a token stream: VulDeePecker's code
gadgets~\cite{li2018vuldeepecker} and SySeVR's syntax- and semantics-based
candidates~\cite{li2022sysevr} sliced programs into data-dependent token
sequences for recurrent classifiers. The Code Property Graph
(CPG)~\cite{yamaguchi2014modeling}, unifying abstract syntax, control flow and
data dependence in one structure, then made graph learning natural, and
Devign~\cite{zhou2019devign} established the template still in use: parse to a
graph, embed node text, propagate, pool, classify. Later work varied the
graph or the propagation --- inter-procedural slices in
DeepWukong~\cite{cheng2021deepwukong}, feature-attention over program
dependence in IVDetect~\cite{li2021ivdetect}, heterogeneous attention
in~\cite{luo2025hagnn}. A third line embeds code with pre-trained language
models, either alone as in LineVul~\cite{fu2022linevul} or fused with a graph
channel~\cite{hin2022linevd, sun2024vullmgnns}.

What is largely invariant across this progression is the node feature: some
embedding of the node's lexical content. Our work isolates that choice. We
hold the graph, the propagation scheme and the training protocol fixed and
vary only what a node carries, which --- to our knowledge --- prior work does
not do as a controlled comparison.

\subsection{Normalisation and Abstraction of Code}
Abstracting away identifiers is long-established practice, but as
preprocessing rather than as an object of study. Clone detection normalises
identifiers and literals so that Type-2 clones become textually
identical~\cite{yamaguchi2014modeling}; sequence-based detectors routinely map
user-defined names to placeholder tokens (\texttt{VAR1}, \texttt{FUN1}) before
training~\cite{li2018vuldeepecker, li2022sysevr}; the Devign corpus itself
ships a normalised variant of every function.

Two properties distinguish what we do. First, existing normalisation is
\emph{anonymising}: it erases the identifier and replaces it with a fresh
symbol, preserving only equality of occurrences. Ours is \emph{semantic}: it
maps a call onto the memory effect the callee is documented to have, so that
\texttt{av\_malloc} and \texttt{kmalloc} become the same symbol as
\texttt{malloc} rather than three distinct placeholders. Second, and more
importantly for a languages venue, we treat the mapping as an analysis with a
specification and a measured error rate, rather than as an unexamined
preprocessing step. We are not aware of prior work in this area that reports
precision and recall for its own abstraction.

The idea of typing operations by their effect is of course not new in
programming languages --- it is the intuition behind effect systems and behavioural
types --- but those require type information and whole-program context that
function-level vulnerability corpora do not provide. Our $\phi$ is deliberately
weaker: syntax-directed, name-based, and total, so it applies to isolated,
non-preprocessed functions, at the cost of the precision quantified in
Section~\ref{sec:rq_analysis}.

\subsection{The Cross-Project Generalisation Gap}
The finding that motivates this paper is now well replicated. Chakraborty et
al.~\cite{chakraborty2022deep} showed that detectors reporting strong
within-project results degrade sharply on realistic, project-disjoint data;
DiverseVul~\cite{chen2023diversevul} demonstrated that simply adding training
projects does not close the gap; PrimeVul~\cite{ding2025primevul} showed that
much apparent progress reflects duplication and label noise, with F1 falling
dramatically under a rigorous split; and~\cite{chakraborty2024realvul} found
the same degradation at repository scale.

Responses have been mostly architectural or adaptational --- domain adaptation
across projects~\cite{liu2022cdvuld, zhang2023cpvd, li2023vulgda}, transfer
with metric fusion~\cite{cai2024csvdtf}. Our result is complementary and, we
argue, prior: with architecture held constant, changing the node
representation alone yields a larger and more consistent improvement than the
architecture comparisons in Section~\ref{sec:rq_generality}. This suggests the
gap is at least partly a representation problem that architectural work has
been indirectly compensating for.

\subsection{Evaluation Methodology in this Area}
The methodological pitfalls we control for are documented in exactly the
literature above, which is why we treat them as requirements rather than
refinements. Duplication across and within corpora inflates results unless
removed before splitting~\cite{ding2025primevul, chen2023diversevul}. Project
mass is heavily concentrated, so a nominally cross-project split can
degenerate into evaluation on a single repository unless large projects are
handled explicitly. Threshold-dependent metrics are unstable when the
calibration and test distributions differ, which they do by construction in
cross-project evaluation. And functions within a repository are not
independent, so confidence intervals computed over functions understate
uncertainty. Section~\ref{sec:protocol} states how each is addressed.

\subsection{Federated and Privacy-Preserving Variants}
A body of work applies federated learning to software engineering tasks,
including defect prediction~\cite{yamamoto2023flr}, code clone detection and
defect prediction jointly~\cite{yang2024federated}, and vulnerability
detection~\cite{zhang2024flvul, hu2024vulfl}, typically with token or metric
features; FedProto~\cite{tan2022fedproto} introduced sharing class prototypes
rather than parameters, and differential privacy for learned models follows
DP-SGD~\cite{abadi2016deep} with R\'enyi accounting~\cite{mironov2019rdp}.
We mention this line of work only to place the present paper against it. This
paper makes no federated or privacy claim, reports no federated experiment, and
evaluates a single centrally trained model throughout. A representation whose
alphabet is fixed by the grammar is a plausible ingredient for such settings,
since it carries no project-specific vocabulary, but we have not tested that
and do not assert it.
